\chapter{Normal subgroups and local quotients}
\label{chap:library:groups}

Throughout this chapter, $G$ and $H$ are finite groups and $\ell$ is a prime.
For a subgroup $Q\leq G$, write $N_G(Q)$ for its normaliser. Since $Q$ is
normal in $N_G(Q)$, the quotient $N_G(Q)/Q$ is defined. We study these
quotients and radical subgroups under isomorphisms and surjective
homomorphisms. They will occur in the definition of a weight in
Chapter~\ref{chap:library:weights}.

\section{The largest normal \texorpdfstring{$\ell$}{ell}-subgroup}

\begin{definition}
The \emph{$\ell$-core} $O_\ell(G)$ is the largest normal $\ell$-subgroup
of $G$. A subgroup $Q\leq G$ is \emph{$\ell$-radical} if
$Q=O_\ell(N_G(Q))$.
\end{definition}

To see that the core exists, take the subgroup generated by all normal
$\ell$-subgroups of $G$. Indeed,
the product of two normal $\ell$-subgroups is normal and has order a power
of $\ell$. Finiteness reduces the generated subgroup to such a product.
It contains every normal $\ell$-subgroup and is characterised by this
property. An automorphism permutes the normal $\ell$-subgroups, so it
preserves the core. More generally, a group isomorphism $\phi:G\to H$ satisfies
$\phi(O_\ell(G))=O_\ell(H)$.

In particular, a radical subgroup is an $\ell$-group. If $\phi:G\to H$ is
an isomorphism, then $\phi(N_G(Q))=N_H(\phi(Q))$, so the equality of cores
shows that $Q$ is radical if and only if $\phi(Q)$ is radical.
\begin{implementation}
\module{ModularRep.PCore} defines the core using Mathlib's subgroup lattice
and its theorem on suprema of normal $\ell$-subgroups.
The basic core results include
\lean{ModularRep.pCore_isPGroup}, \lean{ModularRep.pCore_normal} and
\lean{ModularRep.normal_pSubgroup_le_pCore}.
\module{ModularRep.RadicalSubgroup} defines radicality by mapping the core
of $N_G(Q)$ into $G$ along the normaliser inclusion.
\module{ModularRep.RadicalTransport} proves invariance under isomorphisms.
\end{implementation}

\begin{proposition}[The core lies in every radical subgroup]
\label{lib:groups:core-contained}
Every $\ell$-radical subgroup of $G$ contains $O_\ell(G)$.
\end{proposition}
\begin{proof}
Let $P=O_\ell(G)$ and let $Q$ be radical. Both are $\ell$-groups, and
normality of $P$ makes $PQ$ an $\ell$-group. Suppose that $Q<PQ$.
The normaliser condition in a finite
$\ell$-group gives $x\in N_{PQ}(Q)\setminus Q$. Write $x=yz$, with
$y\in P$ and $z\in Q$. Both $x$ and $z$ normalise $Q$, hence
$y\in P\cap N_G(Q)$. This intersection is a normal $\ell$-subgroup of
$N_G(Q)$ and therefore lies in $O_\ell(N_G(Q))=Q$. Thus $x=yz\in Q$,
a contradiction. It follows that $PQ=Q$ and $P\leq Q$.
\end{proof}
\begin{implementation}
The theorem is \lean{ModularRep.pCore_le_of_isRadicalSubgroup} in
\module{ModularRep.NormalCoreTransport}. It uses Mathlib's nilpotence theorem
for finite $\ell$-groups and its normaliser condition for nilpotent groups.
\end{implementation}

\section{Surjections and normaliser quotients}

\begin{proposition}[Normaliser quotients under a surjection]
\label{lib:groups:normaliser-quotient}
Let $f:G\twoheadrightarrow H$ be surjective, and let $Q\leq G$ contain
$\ker f$. Then $f(N_G(Q))=N_H(f(Q))$, and the map $gQ\mapsto f(g)f(Q)$
induces an isomorphism
\[
 N_G(Q)/Q\;\simeq\;N_H(f(Q))/f(Q).
\]
\end{proposition}
\begin{proof}
The inclusion $f(N_G(Q))\leq N_H(f(Q))$ follows by applying $f$ to a
conjugate of $Q$. Conversely, choose $h\in N_H(f(Q))$ and lift it to
$g\in G$. Since $Q$ contains the kernel, $Q=f^{-1}(f(Q))$. Conjugation
by $g$ therefore preserves $Q$, proving normaliser surjectivity.

Compose this normaliser map with the quotient by $f(Q)$. An element $g$
of $N_G(Q)$ is in its kernel exactly when $f(g)\in f(Q)$, equivalently
when $g\in Q$. The first isomorphism theorem gives the assertion, with
the stated formula on cosets.
\end{proof}

Thus containment of $\ker f$ in $Q$ suffices to identify the normaliser
quotients. To preserve radicality as well, the following result assumes
that $\ker f$ is an $\ell$-group.

\begin{proposition}[Radical subgroups under a surjection]
\label{lib:groups:radical-quotient}
Suppose that $f:G\twoheadrightarrow H$ has $\ell$-group kernel. Then
$f(O_\ell(G))=O_\ell(H)$. For every $Q\leq G$ containing $\ker f$,
$Q$ is $\ell$-radical if and only if $f(Q)$ is $\ell$-radical.
Consequently image and inverse image give inverse correspondences between
the radical subgroups of $G$ and those of $H$.
\end{proposition}
\begin{proof}
The image of the core is a normal $\ell$-subgroup, so
$f(O_\ell(G))\leq O_\ell(H)$. The inverse image of $O_\ell(H)$ is
an extension of two $\ell$-groups and is normal in $G$. It lies in
$O_\ell(G)$, proving equality.

The restriction of $f$ to $N_G(Q)$ is surjective onto $N_H(f(Q))$ by
Proposition~\ref{lib:groups:normaliser-quotient}, and its kernel is an
$\ell$-group. Applying the core equality to this restriction identifies
the images of the two normaliser cores. Both $Q$ and $O_\ell(N_G(Q))$
contain $\ker f$. Image under $f$ is injective on subgroups containing
the kernel, so equality of their images is equivalent to equality of these
subgroups. This proves the radical equivalence.

Finally, $\ker f\leq O_\ell(G)$, so every radical subgroup of $G$
contains $\ker f$ by Proposition~\ref{lib:groups:core-contained}.
The subgroup correspondence theorem now gives the two inverse maps.
\end{proof}
\begin{implementation}
\module{ModularRep.NormalCoreTransport} proves the core image equality,
normaliser surjectivity and the radical equivalence. Its declarations are
\lean{pCore_map_surjective_of_ker_isPGroup},
\lean{normalizerMap_surjective},
\lean{isRadicalSubgroup_iff_map_surjective_of_ker_le} and
\lean{normalizerQuotientEquivOfSurjectiveOfKerLE}, all in
\lean{ModularRep}. These constructions use Mathlib's subgroup
correspondence and quotient homomorphisms.
\end{implementation}

\begin{example}[Quotient by the core]
Set $P=O_\ell(G)$. The quotient $G/P$ has trivial $\ell$-core, and its
radical subgroups correspond to $Q/P$ for radical subgroups $Q$ of $G$.
For each such $Q$, the normaliser quotient is unchanged up to the isomorphism
of Proposition~\ref{lib:groups:normaliser-quotient}. A representation of
this local quotient can therefore be transported without changing its
dimension or irreducibility. Chapter~\ref{chap:library:weights} explains
the resulting construction on weights. If $G$ is itself an $\ell$-group,
Proposition~\ref{lib:groups:core-contained} shows that its only radical
subgroup is $G$.
\end{example}

\section{Central quotients}

The preceding quotient results apply to subgroups containing the kernel.
A central kernel of order prime to $\ell$ gives a different subgroup
correspondence.

\begin{proposition}[Quotients by central $\ell'$-subgroups]
\label{lib:groups:central-prime-to}
Let $f:G\to H$ be a surjective homomorphism of finite groups.
Suppose that $Z=\ker f$ is central and $\ell\nmid|Z|$.
Then $Q\mapsto f(Q)$ is a bijection between the $\ell$-subgroups
of $G$ and those of $H$. For each $\ell$-subgroup $R\leq H$, the unique
$\ell$-subgroup mapping onto $R$ is $O_\ell(f^{-1}(R))$.
For every $\ell$-subgroup $Q\leq G$,
\[
 f^{-1}\bigl(N_H(f(Q))\bigr)=N_G(Q),\qquad
 Q\text{ is radical}\ \Longleftrightarrow\ f(Q)\text{ is radical}.
\]
\end{proposition}
\begin{proof}
Let $R\leq H$ be an $\ell$-subgroup. A Sylow $\ell$-subgroup $Q$ of
$f^{-1}(R)$ has order $|R|$, since the kernel has order prime to
$\ell$. Since $Q\cap Z=1$, its image is $R$, and
$f^{-1}(R)=QZ$. Centrality of $Z$ makes $Q$ normal in this
inverse image, so $Q$ lies in its $\ell$-core. That core also
maps onto $R$.

The restriction of $f$ to any $\ell$-subgroup is injective,
since its kernel has order dividing both an $\ell$-power and
$|Z|$. If two $\ell$-subgroups have image $R$, they both lie
in the $\ell$-core of $f^{-1}(R)$. Injectivity on that core
then shows that they are equal. This proves the correspondence
and the formula for its inverse.

An element $g\in G$ normalises $Q$ precisely when $f(g)$
normalises $f(Q)$. For the reverse implication, use uniqueness
of the lift on $Q$ and $gQg^{-1}$. This proves the normaliser
formula and surjectivity of $N_G(Q)\to N_H(f(Q))$.

Uniqueness also implies that the lift of a normal
$\ell$-subgroup is normal. Consequently
$f(O_\ell(G))=O_\ell(H)$. Apply the same argument to the
surjection of normalisers. The image of $O_\ell(N_G(Q))$ is
$O_\ell(N_H(f(Q)))$. Injectivity on $\ell$-subgroups now
proves the radicality equivalence.
\end{proof}

In particular, the induced map on local quotients is
\[
 N_G(Q)/Q\longrightarrow N_H(f(Q))/f(Q).
\]
Its kernel is $QZ/Q\simeq Z$, since $Q\cap Z=1$. Quotienting by this
kernel gives
\[
 N_G(Q)/(QZ)\cong N_H(f(Q))/f(Q).
\]
The two local quotient orders therefore have equal $\ell$-parts.
Inflation along this map gives characters of $N_G(Q)/Q$ that are
trivial on $QZ/Q$.

Taking $f$ to be the quotient map by a central $\ell'$-subgroup $Z$
gives, for every $\ell$-radical subgroup $Q$,
\[
 QZ/Z\text{ is radical in }G/Z,
 \qquad N_{G/Z}(QZ/Z)=N_G(Q)/Z.
\]
These conclusions also follow from
\cite[Lemma~2.3(c) and its proof]{NavarroTiep2011}.
Proposition~\ref{lib:groups:normaliser-quotient} does not apply directly
to $Q$, because $Z$ need not be contained in $Q$.

\begin{implementation}
The subgroup correspondence is proved in
\module{ModularRep.PaperProofs.SporadicFi24P3Definition44NamedCarrierCentralPrimeToPSubgroups},
including \lean{existsUnique_pSubgroup_lift} and
\lean{pSubgroupLift_map}. The module whose name ends in
\lean{CentralPrimeToRadicals} proves
\lean{normalizer_comap_of_central_primeTo} and
\lean{radical_iff_map_of_central_primeTo}. The theorem
\lean{fixedCentralQuotientSource} gives the quotient specialisation
as the structure
\lean{ModularRep.NavarroTiep23cFixedCentralQuotientSource}, defined in
\module{ModularRep.PaperProofs.NavarroTiep23cFixedCentralQuotientSource}.
\module{ModularRep.PaperProofs.CentralEllPrimeWeightLocalQuotient}
defines the homomorphism \lean{qW} on normaliser quotients and computes
its kernel. Its theorems on surjectivity and the $\ell$-parts of the
quotient orders are stated for radical $Q$ and use the quotient
specialisation above. The general subgroup correspondence also allows
the trivial subgroup.
\end{implementation}

\section{Cyclic quotients and Sylow subgroups}

\begin{proposition}[Cyclic quotients by the core]
\label{lib:groups:hypoelementary}
If $G/O_\ell(G)$ is cyclic, then $H/O_\ell(H)$ is cyclic for every
subgroup $H\leq G$. If $G$ is cyclic, $G/O_\ell(G)$ has order prime
to $\ell$.
\end{proposition}
\begin{proof}
Let $I=H\cap O_\ell(G)$. This is a normal $\ell$-subgroup of $H$,
so $I\leq O_\ell(H)$. Restriction of the quotient map embeds $H/I$
in $G/O_\ell(G)$. Thus $H/I$ is cyclic, as is its quotient
$H/O_\ell(H)$. For the second assertion, the unique Sylow $\ell$-subgroup
of a finite cyclic group is normal and equals its $\ell$-core. Its index
is prime to $\ell$.
\end{proof}

For a finite group the condition that $G/O_\ell(G)$ be cyclic is
equivalent to $\ell$-hypoelementarity, that is, the existence of a normal
$\ell$-subgroup with cyclic quotient of order prime to $\ell$.
For this terminology, see, for example,
\cite[Definition~5.5.3]{Benson1998}. Indeed, the inverse image of the
Sylow $\ell$-subgroup of a cyclic quotient would be a normal
$\ell$-subgroup of $G$. Maximality of the core forces that quotient
Sylow subgroup to be trivial. Thus the cyclic quotient has $\ell'$-order.
The prime defining hypoelementarity is specified independently of any
coefficient field.
\begin{implementation}
\lean{ModularRep.subgroup_quotient_pCore_isCyclic} and
\lean{ModularRep.not_dvd_card_quotient_pCore_of_isCyclic}.
The latter belongs to \module{ModularRep.CyclicPCore}.
\end{implementation}

\begin{proposition}[A criterion for cyclic Sylow subgroups]
\label{lib:groups:cyclic-sylow}
Let $T\leq G$ be cyclic. If $|G|_\ell$ divides $|T|$, then $G$ has
a cyclic Sylow $\ell$-subgroup, and every $\ell$-subgroup of $G$ is cyclic.
Here $|G|_\ell$ denotes the largest power of $\ell$ dividing $|G|$.
\end{proposition}
\begin{proof}
The cyclic group $T$ contains a subgroup of order $|G|_\ell$. It is a
Sylow subgroup of $G$ by its order. Every $\ell$-subgroup of $G$ is
contained in a Sylow subgroup, and all Sylow subgroups are conjugate.
Hence each is contained in a cyclic group and is cyclic.
\end{proof}
\begin{implementation}
\module{ModularRep.CyclicSylow}, in particular
\lean{exists_cyclic_sylow_of_full_part_dvd_card_cyclic_subgroup} and
\lean{isCyclic_of_isPGroup_of_full_part_dvd_card_cyclic_subgroup}
in that module's namespace. They combine Mathlib's Sylow existence,
containment and conjugacy theorems with transport of cyclicity.
\end{implementation}

\begin{example}[Applying the order criterion]
Suppose $G$ has a cyclic subgroup $T$ and an order
factorisation $|G|=ab$ with $\ell\nmid a$ and $b\mid |T|$.
Then $|G|_\ell\mid b$, so the preceding proposition gives cyclicity of
every $\ell$-subgroup of $G$. This divisibility is proved by
\lean{ModularRep.CyclicSylow.ordProj_mul_dvd_right_of_not_dvd_left}.
\end{example}
